\section{Матричный алгоритм достижимости с ограничениями в виде конъюнктивных и булевых языков}
	\label{sec:ConjunctiveBooleanReachability}
	\tikzsetfigurename{BeyondCFL_ConjBool_}

Конъюнктивные и булевы грамматики позволяют задавать ограничения на пути, которые не выражаются контекстно-свободными грамматиками. Это полезно, например, в статическом анализе программ, где линейные конъюнктивные грамматики применяются для анализа зависимостей данных~\sidecite{10.1145/3093333.3009848}. Однако увеличение выразительной силы приводит к принципиальным ограничениям: задача поиска путей с ограничениями в виде конъюнктивных и булевых языков в реляционной семантике запросов в общем случае неразрешима~\sidecite{hellingsRelational}. Поэтому далее мы рассматриваем алгоритмы, которые строят аппроксимацию ответа сверху: если между двумя вершинами действительно существует путь, удовлетворяющий ограничению, то соответствующая пара будет найдена, но в результате могут появиться лишние пары.

Для помеченного графа $D = (V,E,\Sigma)$ и грамматики $G = (\Sigma,N,P,S)$ удобно рассматривать не одно отношение для стартового нетерминала, а семейство отношений для всех нетерминалов. Для каждого $A \in N$ определим
\[
    R_A = \{(i,j) \in V \times V \mid \exists \pi: i \pi j \land l(\pi) \in L(A)\}.
\]
Здесь $i \pi j$ означает, что путь $\pi$ начинается в вершине $i$ и заканчивается в вершине $j$, а $l(\pi)$~--- слово, полученное конкатенацией меток рёбер пути. В реляционной семантике результатом запроса для стартового нетерминала $S$ является отношение $R_S$. Алгоритмы ниже строят матрицы $T$, элементы которых являются подмножествами $N$; условие $A \in T_{i,j}$ интерпретируется как найденное свидетельство того, что пара $(i,j)$ принадлежит аппроксимации отношения $R_A$.

\subsection{Алгоритм для конъюнктивных грамматик}
	\label{subsec:ConjunctiveReachabilityAlgorithm}

Пусть конъюнктивная грамматика $G = (\Sigma,N,P,S)$ приведена к бинарной нормальной форме в смысле теоремы~\ref{thm:BinaryNormalForm}. Для матриц $X$ и $Y$ подходящих размеров, элементы которых являются подмножествами $N$, определим конъюнктивное произведение $X \circ Y = Z$ следующим образом. Сначала для каждой пары индексов вычисляется множество пар нетерминалов
\[
    M_{i,j} = \bigcup_k X_{i,k} \times Y_{k,j}.
\]
После этого
\[
    Z_{i,j} = \{A \mid \exists (A \rightarrow B_1 C_1 \& \ldots \& B_m C_m) \in P: \forall t, 1 \leq t \leq m, (B_t,C_t) \in M_{i,j}\}.
\]
Иными словами, нетерминал $A$ добавляется в ячейку $(i,j)$, если для каждого конъюнкта правила $A$ в этой ячейке найдено разбиение пути на две части, соответствующие нетерминалам этого конъюнкта. При этом разные конъюнкты могут опираться на разные пути из $i$ в $j$, что и порождает аппроксимацию сверху.

Конъюнктивное транзитивное замыкание матрицы $T$ определяется итеративно:
\[
    T^{(1)} = T, \qquad T^{(r)} = T^{(r-1)} \cup (T^{(r-1)} \circ T^{(r-1)}), \quad r \geq 2.
\]
Предел этой последовательности обозначим $T^{\mathit{conj}}$. Так как множество возможных значений конечно, вычисление можно выполнять до неподвижной точки.

\begin{algorithm}[H]
    \floatname{algorithm}{Листинг}
    \footnotesize
    \caption{Матричный алгоритм для конъюнктивных грамматик}
    \label{lst:conj_matrix_reachability}
    \KwData{$D = (V,E,\Sigma)$~--- помеченный граф; $G = (\Sigma,N,P,S)$~--- конъюнктивная грамматика в бинарной нормальной форме}
    \KwResult{Матрица $T$, задающая аппроксимации отношений $R_A$ для всех $A \in N$}
    $n \gets |V|$\;
    $T \gets$ матрица размера $n \times n$, заполненная $\varnothing$\;
    \ForEach{$(i,a,j) \in E$}{
        $T_{i,j} \gets T_{i,j} \cup \{A \mid (A \rightarrow a) \in P\}$\;
    }
    \While{$T$ изменяется}{
        $T \gets T \cup (T \circ T)$\;
    }
    \Return{$T$}\;
\end{algorithm}

Алгоритм на листинге~\ref{lst:conj_matrix_reachability} был предложен Азимовым и Григорьевым как обобщение матричного алгоритма КС-достижимости~\sidecite{Azimov:2018:CPQ:3210259.3210264} на конъюнктивные грамматики~\sidecite{565CECD7E8F5C6063935B41DB41797AA37D53B04}. Инициализация добавляет в матрицу факты, соответствующие терминальным правилам и рёбрам графа. Затем итерации замыкания добавляют нетерминалы, для которых выполнены все конъюнкты некоторого бинарного правила.

\begin{example}[Аппроксимация для конъюнктивной грамматики]
    Рассмотрим грамматику $G = (\Sigma,N,P,S)$, где $N = \{S,A,B,C,D\}$, $\Sigma = \{a,b,c\}$, а правила имеют вид
    \[
    \begin{array}{rclcrcl}
        S &\rightarrow& AB \& DC, && A &\rightarrow& a, \\
        B &\rightarrow& BC \mid b, && C &\rightarrow& c, \\
        D &\rightarrow& AD \mid b. &&&&
    \end{array}
    \]
    Конъюнкт $AB$ порождает язык $\{abc^*\}$, а конъюнкт $DC$~--- язык $\{a^*bc\}$. Поэтому $L(S)=\{abc\}$.

    Пусть граф $D$ имеет рёбра
    \[
        (0,a,1), (1,b,2), (2,c,3), (3,c,4), (1,a,5), (5,b,6), (6,c,4).
    \]
    После инициализации матрица $T_0$ содержит $A$ в ячейках $(0,1)$ и $(1,5)$, $B$ и $D$ в ячейках $(1,2)$ и $(5,6)$, а $C$ в ячейках $(2,3)$, $(3,4)$ и $(6,4)$. На следующих итерациях появляются, например, $D \in T_{0,2}$ и $B \in T_{1,3}$, а затем $S \in T_{0,3}$, что соответствует пути с меткой $abc$.

    В неподвижной точке аппроксимации отношений имеют вид
    \[
    \begin{array}{rcl}
        R'_S &=& \{(0,3),(0,4),(1,4)\},\\
        R'_A &=& \{(0,1),(1,5)\},\\
        R'_B &=& \{(1,2),(1,3),(1,4),(5,4),(5,6)\},\\
        R'_C &=& \{(2,3),(3,4),(6,4)\},\\
        R'_D &=& \{(0,2),(0,6),(1,2),(1,6),(5,6)\}.
    \end{array}
    \]
    Пара $(0,4)$ является лишней для $R_S$: между вершинами $0$ и $4$ есть путь с меткой $abcc$, соответствующий конъюнкту $AB$, и путь с меткой $aabc$, соответствующий конъюнкту $DC$, но нет пути с меткой $abc$. Алгоритм не различает, что конъюнкты были подтверждены разными путями между той же парой вершин, поэтому добавляет $S$ в $T_{0,4}$.
\end{example}

\begin{theorem}[Завершаемость]
    Для любого конечного помеченного графа $D = (V,E,\Sigma)$ и любой конечной конъюнктивной грамматики $G = (\Sigma,N,P,S)$ алгоритм~\ref{lst:conj_matrix_reachability} завершается.
\end{theorem}

\begin{proof}
    На каждой итерации алгоритм только добавляет нетерминалы в ячейки матрицы $T$ и никогда их не удаляет. Всего существует не более $|V|^2|N|$ таких добавлений. Следовательно, после конечного числа итераций матрица перестанет изменяться.
\end{proof}

\begin{theorem}[Полнота аппроксимации]
    Пусть $T$~--- матрица, полученная алгоритмом~\ref{lst:conj_matrix_reachability}. Тогда для любых $i,j \in V$ и $A \in N$, если $(i,j) \in R_A$, то $A \in T_{i,j}$.
\end{theorem}

\begin{proof}
    Доказательство проводится индукцией по высоте дерева вывода слова $l(\pi)$ для пути $i \pi j$. Если высота равна единице, то вывод использует терминальное правило $A \rightarrow a$, а соответствующий факт добавляется при инициализации. Для индукционного перехода рассмотрим правило $A \rightarrow B_1 C_1 \& \ldots \& B_m C_m$ в корне дерева вывода. Для каждого конъюнкта существует разбиение соответствующего пути на части, выводимые из $B_t$ и $C_t$. По индукционному предположению эти нетерминалы уже будут добавлены в соответствующие ячейки на некоторой предыдущей итерации. Тогда конъюнктивное произведение добавит $A$ в ячейку $(i,j)$.
\end{proof}

Если размер грамматики считается константой, то одна итерация вычисления $T \cup (T \circ T)$ сводится к константному числу умножений, объединений и пересечений булевых матриц размера $|V| \times |V|$. Обозначим время этих операций через $\BMM{|V|}$, $\BMU{|V|}$ и $\BMI{|V|}$ соответственно. Так как число итераций ограничено $|V|^2|N|$, при фиксированной грамматике получаем оценку
\[
    O(|V|^2(\BMM{|V|} + \BMU{|V|} + \BMI{|V|})).
\]

\subsection{Алгоритм для булевых грамматик на ациклических графах}
	\label{subsec:BooleanDAGReachabilityAlgorithm}

Конъюнктивный алгоритм не переносится напрямую на булевы грамматики: отрицательные конъюнкты требуют знать не только наличие, но и отсутствие вывода для той же строки. Для произвольных графов это особенно проблематично из-за циклов и потенциально бесконечного множества путей. В работе Шеметовой и Григорьева предложен матричный алгоритм для случая, когда входной граф является ориентированным ациклическим графом~\sidecite{Shemetova2019}. Он строит аппроксимацию сверху для булевой грамматики в двоичной нормальной форме.

Пусть вершины DAG упорядочены топологически, так что каждое ребро ведёт из вершины с меньшим номером в вершину с большим номером. Тогда все существенные значения лежат в верхнетреугольной части матриц. Алгоритм использует две таблицы. Первая таблица $T$ размера $n \times n$ хранит множества нетерминалов: $A \in T_{i,j}$ означает, что пара $(i,j)$ добавлена в аппроксимацию $R_A$. Вторая таблица $M$ хранит множества пар нетерминалов из $N \times N$ и аккумулирует информацию о разбиениях путей.

Для булевой грамматики в двоичной нормальной форме определим функцию $f:2^{N \times N}\rightarrow 2^N$. Нетерминал $A$ принадлежит $f(M)$, если существует правило
\[
    A \rightarrow B_1 C_1 \& \ldots \& B_m C_m \& \neg D_1 E_1 \& \ldots \& \neg D_k E_k
\]
и существует подмножество $M' \subseteq M$, такое что все положительные пары $(B_t,C_t)$ принадлежат $M'$, а все отрицательные пары $(D_t,E_t)$ не принадлежат $M'$. Рассмотрение подмножеств необходимо для полноты аппроксимации: разные пары в $M$ могли быть получены из разных путей между одними и теми же вершинами, и запрет пары для одного пути не должен удалять корректное свидетельство, полученное по другому пути.

\begin{algorithm}[H]
    \floatname{algorithm}{Листинг}
    \footnotesize
    \caption{Матричный алгоритм для булевых грамматик на DAG}
    \label{lst:boolean_dag_matrix_reachability}
    \KwData{$D = (V,E,\Sigma)$~--- помеченный DAG; $G = (\Sigma,N,P,S)$~--- булева грамматика в двоичной нормальной форме}
    \KwResult{Матрица $T$, задающая аппроксимации отношений $R_A$ для всех $A \in N$}
    $D \gets$ результат топологической сортировки $D$\;
    $n \gets |V|$\;
    $T,M \gets$ пустые матрицы размера $n \times n$\;
    \ForEach{$(i,a,j) \in E$}{
        $T_{i,j} \gets T_{i,j} \cup \{A \mid (A \rightarrow a) \in P\}$\;
    }
    $compute(0,n)$\;
    \Return{$T$}\;
\end{algorithm}

Основная работа выполняется рекурсивными процедурами $compute$ и $complete$. Процедура $compute(l,m)$ вычисляет значения $T_{i,j}$ для всех $l \leq i < j < m$. Она разбивает диапазон пополам, рекурсивно обрабатывает левую и правую половины, а затем вызывает $complete$ для прямоугольной области, соединяющей левую половину с правой. Процедура $complete(l,m,l',m')$ обрабатывает прямоугольную подматрицу, соответствующую путям, начало которых лежит в диапазоне $[l,m)$, а конец~--- в диапазоне $[l',m')$. В базовом случае, когда прямоугольник состоит из одной ячейки, значение $T_{l,l'}$ обновляется как $f(M_{l,l'})$. В рекурсивном случае процедура разбивает прямоугольник на меньшие области и обновляет соответствующие блоки $M$ произведениями блоков $T$.

\mytodo{Перерисовать схему расположения подматриц для процедуры compute по статье Шеметовой и Григорьева~\cite{Shemetova2019}.}

\mytodo{Перерисовать схему расположения подматриц для процедуры complete по статье Шеметовой и Григорьева~\cite{Shemetova2019}.}

Смысл рекурсивной схемы такой же, как в матричном алгоритме Охотина для булевых грамматик~\sidecite{Okhotin:2014:PMM:2565359.2565379}: значения для более коротких фрагментов вычисляются раньше, после чего произведение подматриц переносит в $M$ все пары нетерминалов, которые могут быть получены при разбиении пути промежуточной вершиной. Топологический порядок гарантирует, что такое динамическое программирование по верхнетреугольной части матрицы корректно для DAG.

\begin{example}[Аппроксимация для булевой грамматики на DAG]
    В статье~\sidecite{Shemetova2019} работа алгоритма демонстрируется на булевой грамматике, содержащей правило с отрицательным конъюнктом
    \[
        S \rightarrow DC \& \neg AB.
    \]
    После топологической сортировки графа алгоритм инициализирует таблицу $T$ по терминальным правилам, а затем последовательно выполняет рекурсивные вызовы $compute$ и $complete$. Для примера из статьи дерево вызовов начинается с $compute(0,8)$, затем обрабатываются $compute(0,4)$ и $compute(4,8)$, после чего вызывается $complete(0,4,4,8)$ для области, соединяющей две половины таблицы.

    \mytodo{Перерисовать дерево рекурсивных вызовов compute/complete из примера алгоритма для булевых грамматик по статье Шеметовой и Григорьева~\cite{Shemetova2019}.}

    При обработке одной из ячеек в $M_{4,7}$ оказывается множество пар, содержащее как $(D,C)$, так и $(A,B)$. Если бы функция $f$ рассматривала всё множество целиком, правило $S \rightarrow DC \& \neg AB$ не могло бы быть применено из-за наличия пары $(A,B)$. Однако пары могли соответствовать разным путям между вершинами $4$ и $7$: один путь подтверждает положительный конъюнкт $DC$, а другой даёт пару, запрещённую отрицательным конъюнктом. Поэтому $f$ рассматривает подмножества $M_{4,7}$ и может добавить $S$ по подмножеству, где есть $(D,C)$ и нет $(A,B)$.

    Это обеспечивает полноту аппроксимации, но может добавить лишние ответы. В том же примере в итоговой таблице появляется $S \in T_{4,7}$, хотя путь с нужной меткой для $S$ отсутствует: положительный и отрицательный факты были получены из разных путей между одной парой вершин.

    \mytodo{Восстановить таблицы состояний $T$ из примера статьи Шеметовой и Григорьева в виде TeX/TikZ-рисунков; не использовать изображения, извлечённые из PDF, как финальные иллюстрации.}
\end{example}

\begin{theorem}[Полнота аппроксимации для DAG]
    Пусть $D = (V,E,\Sigma)$~--- помеченный DAG, $G = (\Sigma,N,P,S)$~--- булева грамматика в двоичной нормальной форме, а $T$~--- матрица, построенная алгоритмом~\ref{lst:boolean_dag_matrix_reachability}. Тогда для любых $i,j \in V$ и $A \in N$, если $(i,j) \in R_A$, то $A \in T_{i,j}$.
\end{theorem}

\begin{proof}
    Доказательство проводится индукцией по высоте дерева вывода для слова $l(\pi)$, где $i \pi j$. В базе индукции путь состоит из одного ребра, а соответствующий нетерминал добавляется в $T$ при инициализации. Для индукционного перехода рассмотрим правило булевой грамматики в корне дерева вывода. Так как граф топологически упорядочен, все пути, соответствующие поддеревьям, лежат в подинтервалах $[i,j]$ и вычисляются раньше той ячейки, которая соответствует всему пути. По индукционному предположению нетерминалы для положительных конъюнктов будут присутствовать в соответствующих ячейках $T$, а произведения подматриц добавят нужные пары в $M_{i,j}$. Функция $f$ рассматривает подходящее подмножество $M_{i,j}$, соответствующее данному пути, поэтому отрицательные конъюнкты проверяются относительно этого подмножества, а не относительно всех пар, накопленных из всех путей между $i$ и $j$. Следовательно, $A$ будет добавлен в $T_{i,j}$.
\end{proof}

Сложность алгоритма совпадает с оценкой матричного алгоритма Охотина для булевых грамматик на линейном входе с заменой длины строки на число вершин DAG. Если $\BMM{n}$~--- время умножения булевых матриц размера $n \times n$, то алгоритм работает за
\[
    O(|G|\BMM{|V|}\log |V|).
\]
Ограничение на ацикличность существенно: оно позволяет использовать топологический порядок и вычислять значения динамически по верхнетреугольной матрице. Вопрос о сохранении такой же асимптотики для произвольных графов остаётся открытым~\sidecite{Shemetova2019}.
